\documentclass[border = 1mm]{standalone}
% Figure 2
% Geometry of the n-th inclined plane
%
% Source:
% R. Sánchez-Martínez and E. Heredia-Muñoz
% Pendulum Period from a Sequence of Inclined Planes
%
% Shared for educational and research purposes.
%
% Figure source archived at:
% github.com/ro-smtz/research
%
% Generated with TikZ and compiled preferable with XeLaTeX.

% -- Preamble
% input
\usepackage[utf8]{inputenc}
\usepackage[T1]{fontenc}
% font (the two packages require XeLaTeX)
\usepackage{fontspec}
\usepackage[default]{fontsetup}
% TikZ
\usepackage{tikz}
\usetikzlibrary{
	math,calc,
	angles,quotes,
	intersections,
	arrows.meta,bending
}

% -- Main document
\begin{document}

% Figure illustrating the geometric construction of the n-th
% inclined plane used to approximate the pendulum trajectory.
% Note that the triangles $\triangle \xi_n x_n \xi_n'$ and
% $\triangle \bar{x}_n O \xi_n'$ are similar, since both are
% right triangles sharing the angle at $\xi_n'$. This allows
% us to write the inclination angle of the $n$-th plane in 
% terms of the angle from the vertical as $\theta_n = \phi_n
% + \frac{1}{2}\delta\phi$.
\begin{tikzpicture}[
	line join = round,
	line cap = round,
	line width = 0.6pt,
	>/.tip = {Straight Barb[width = 1.5mm, length = 1mm, line cap = round]},
]
	
	% Geometric parameters:
	% l              : pendulum length
	% ph{x0}         : initial angular displacement
	% ph{xa}, ph{xb} : consecutive discretization points x_{n-1} and x_n
	% dph            : angular increment \delta\phi
	% th             : local inclination angle \theta_n of the n-th plane
	\tikzmath{
		\l = 5.2;
		\ph{x0} = 56;
		\ph{xa} = 43; \ph{xb} = 30;
		\dph = \ph{xa} - \ph{xb};
		\th = \ph{xb} + 0.5*\dph;
		\dth = 10;
	}
	
	% Coordinates defining the circular arc and the n-th inclined plane
	\coordinate (O) at (0,0);
	
	\coordinate (x0) at (-90+\ph{x0}:\l);
	\coordinate (xa) at (-90+\ph{xa}:\l);
	\coordinate (xb) at (-90+\ph{xb}:\l);
	\coordinate (xN) at (0,-\l);
	
	\coordinate (xm) at ($(xa)!0.5!(xb)$);
	
	% arc
	\draw[gray!40] (O)--(\l,0) arc (0:-90:\l) --cycle;
	
	% Auxiliary construction used to locate \xi'_n,
	% the intersection between the extension of the n-th plane
	% and the vertical line through x_N.	
	\begin{scope}
		\clip (O) rectangle ++(\l,-1.3*\l);
		
		\path[name path = Oc] (xN) -- ++(0,-0.5*\l);
		\path[name path = ab] (xa)--($(xa)!6!(xb)$);
		\path[name intersections = {of = ab and Oc, by = xc}];
	\end{scope}
	
	\coordinate (xi) at (xb-|O);
	
	% Drawing the initial angle \Delta\phi
	\draw[<->, draw = gray!25]
		(-90:0.5*\l) arc (-90:-90+\ph{x0}:0.5*\l)
		node (tmp) [
			inner sep = 1mm,
			pos = 0.3, text = gray, fill = white,
			shift = {(2pt,0pt)},
		]
		{\phantom{$\Delta\phi$}};
		% creates an empty node first to interrupt the angle stroke
	\node[
		text = gray,
		xshift = 0.5mm,
		yshift = 0.2mm,
	] at (tmp) {$\Delta\phi$};
	
	% Drawing the inclination angle \theta_n
	\draw[<->, draw = gray!25]
		(-90:0.67*\l) arc (-90:-90+\th:0.67*\l)
		node (tmp) [
			inner sep = 1mm,
			pos = 0.4, text = gray, fill = white,
			shift = {(2pt,0pt)},
		]
		{\phantom{$\theta_n$}};
		% creates an empty node first to interrupt the angle stroke
	\node[
		text = gray,
		xshift = 0.5mm,
		yshift = 0.2mm,
	] at (tmp) {$\theta_n$};
	
	% Dimension marker indicating the discretization interval \delta x
	\path
		(xb) to (xa)
		to ([turn]-90:0.85) coordinate (xa1)
		to ([turn]0:0.25) coordinate (xa2);
		
	\path
		(xa) to (xb)
		to ([turn]90:0.85) coordinate (xb1)
		to ([turn]0:0.25) coordinate (xb2);
	
	\draw[draw = gray!50]
		(xa1) to coordinate (m1) (xa2)
		(xb1) to coordinate (m2) (xb2)
		%
		(m1) to
			node (n) [sloped, inner sep = 1mm, fill = white] 
			{\phantom{$\delta x$}}
		(m2);
	
	\node[shift = {(0.4mm,0.2mm)}]
		at (n) {$\delta x$};
	
	% Secondary construction lines shown in light gray
	\begin{scope}[gray!40]
		\draw (O) to (x0);
		\draw (O) edge ++(0,3mm) edge ++ (-3mm,0);
		\draw (O) to (xm);
		\draw (xN) edge ++(0.38*\l,0) edge (xc);
		
		\draw[|-|, gray!40]
			([xshift = -7mm]O)
			to node[black, fill = white, pos = 0.6] {$l$}
			([xshift = -7mm]xN);
			
		\draw (xc) edge ++(0,-3mm) edge ++(5mm,0);
		\draw (xN) to ++(-3mm,0);
	\end{scope}
			
	\coordinate (xbb) at (xb-|xa);
			
	% Secondary construction lines shown in light gray
	\pic[
		draw, "$\phi_n$",
		angle radius = 8mm,
		angle eccentricity = 1.4,
	] {angle = xN--O--xb};
	
	\pic[
		draw, "$\delta\phi$",
		angle radius = 1.9cm, 
		angle eccentricity = 1.15,
	] {angle = xb--O--xa};
	
	\pic[draw, angle radius = 5mm] {angle = xa--xc--O};
	
	\draw[shorten >= -1mm, shorten <= 4mm]
		(xc) to ++({0.5*\th + 45}:5mm);
	
	% theta_n in inclined plane
	\pic[
		draw,
		angle radius = 6.4mm,
		"$\theta_n$",
		angle eccentricity = 1.3,
	] {angle = xi--xb--xc};	
	
	\draw[gray!50, shorten >= -3mm] (xb) to (xc);
	
	\pic[
		draw,
		"\small$\theta_n$" {xshift = 0mm, yshift = -0.4mm},
		angle radius = 5mm, angle eccentricity = 1.55,
	] {angle = xbb--xb--xa};
	
	\pic[draw, scale = 0.4] {right angle = O--xm--xb};
	\pic[draw, scale = 0.5] {right angle = xc--xi--xb};
	
	% Principal geometric elements of the n-th inclined plane
	\draw
		(xa) to node[above right = -0.5mm, pos = 0.4] {$l$}
		(O)  to (xb);
	\draw[line width = 1pt] (xa) to (xb);
	\draw[gray] (xb) -| (xa);
	
	\draw[gray] (xb) to (xi) to (xc) to cycle;
	
	% Filled and open markers distinguishing primary and auxiliary points
	\foreach \p in {x0,xa,xb,xN}
		\fill (\p) circle (1.3pt);
		
	\foreach \p in {O,xi,xc}
		\draw[semithick, fill = white] (\p) circle (1.2pt);
	
	\draw[thick, fill = white] (xm) circle (1.2pt);
		
	% labels
	\node[above, xshift = -2.5mm] at (O) {$O$};
	
	\node[right, yshift = -1.5mm] at (x0) {$x_0$};
	\node[shift = {(5mm,-1mm)}] at (xa) {$x_{n-1}$};
	\node[shift = {(0.5mm,3.5mm)}] at (xm) {$\overbar{x}_n$};
	\node[shift = {(1.0mm,-2.8mm)}] at (xb) {$x_n$};
	\node[left = -2pt, yshift = -2.3mm] at (xN) {$x_N$};
	
	\node[above right] at (xi) {$\xi_n$};
	\node[below right] at (xc) {$\xi'_n$};
	
\end{tikzpicture}
\end{document}